\subsection{一般の $q$ における忠実な差分模型と全候補pencil}
\label{subsec:quantum-triangle-general}
$q\ne0,\pm1$ とする。$V_q$ を $yz-qzy=x$ cyclic、$B_q$ をその中央Rees環とする。
\[
 \Omega_q=xyz+\frac{x^2+q^2y^2+z^2}{q^2-1},\qquad
 U_{q,D}=V_q/(\Omega_q+D)
\]
と置く。

\begin{theorem}[全 $q,D$ の整域性と差分表示]
\label{thm:quantum-triangle-faithful-difference}
$B_q$ はnoetherian quadratic Koszul CY AS4であり、$\Omega_q$ は中心元である。
全 $D$ で $U_{q,D}$ と $R_{q,D}=B_q/(\Omega_{q,h}+Dt^3)$ は整域となる。
$U_{q,D}$ の基底は $x^iy^jz^k$、$i,j,k\ge0$、$ijk=0$ である。

$\rho^2=q$、$\kappa=-\rho/(q^2-1)$ とし、
$\mathcal T g(u)=g(qu)\mathcal T$ を満たす差分作用素環を使う。
\[
 b_D(u)=\frac{q(u^2+1)(u^2+q^2)+u^2(q^2-1)^3D}
                  {(q^2-1)^2(u^2-q^2)(u^2-1)}
\]
と置くと、忠実表示は
\[
 \begin{aligned}
 x&=\kappa(u+u^{-1}),\\
 y&=\mathcal T+b_D(u)\mathcal T^{-1},\\
 z&=\rho\bigl(u\mathcal T+u^{-1}b_D(u)\mathcal T^{-1}\bigr)
 \end{aligned}
\]
で与えられる。
\end{theorem}
\begin{proof}
$V_q$ の唯一のPBW重なり $zyx$ は、どちらから還元しても
$q^{-1}(xyz-x^2+y^2-z^2)$ となる。
従って先頭環は量子多項式環 $S_q$ である。
中央ReesのAS4性はcyclic tensorの対称PBW条件と正則中心拡大から従う。
$\mu_{S_q}=1$ で、$Z(S_q)_1=0$ だから $\mu_{B_q}(t)$ は $t$ の倍数である。
homological determinant1と正規商公式でその倍数は1となり、
残る $x_i\mapsto x_i+v_it$ の平行移動も、二次関係の $tX$ 係数と $q\ne1$ から零となる。
従って $B_q$ はCYである。

表示した差分作用素を代入すると三関係が成立し、$\Omega_q$ の像は厳密に $-D$ となる。
まず $D$ を独立な変数としてこの写像の忠実性を示す。
$\Omega_q$ の最高項は $xyz$ だから
$\Omega_q^l x^iy^jz^k$、$ijk=0$、$l\ge0$ は $V_q$ のfiltered基底となる。
ここではまだ中心性を使っていない。

$n=j+k$ ごとに最高 $\mathcal T$ 係数を取る。
$n>0$ では、非零定数倍を除き、その線型結合は
\[
 F(x(u))+u^nG(x(u))+\sum_{k=1}^{n-1}c_ku^k
\]
の形になる。
これが零なら $u\mapsto u^{-1}$ の式との差を $u-u^{-1}$ で割る。
$(u^k-u^{-k})/(u-u^{-1})$ は $x(u)$ の次数 $k-1$ 多項式なので、
次数 $n-1$ の項から $G=0$、続いて全 $c_k=0$、最後に $F=0$ を得る。
$n=0$ は $x(u)$ の多項式の独立性である。
$\Omega_q^l$ の像 $(-D)^l$ も独立だから、$V_q$ の写像は忠実である。
従って $\Omega_q$ は中心となる。
さらにこの基底は $k[\Omega_q]$ 上の自由基底を与え、
各固定 $D$ の商でも上記最高係数の独立性は変わらない。
従って各 $U_{q,D}$ も差分作用素整域へ忠実に埋まり、Rees環 $R_{q,D}$ も整域となる。
\end{proof}

\begin{proposition}[全候補incidenceと有限位数のChebyshev中心元]
\label{prop:quantum-triangle-all-candidates}
$q$ が非捩れなら $Z(V_q)=k[\Omega_q]$、$Z(B_q)=k[t,\Omega_{q,h}]$ である。
$q$ が正確に位数 $N\ge3$ なら、
\[
 c_0(X)=2,\quad c_1(X)=X,\quad
 c_n(X)=Xc_{n-1}(X)-\kappa^2c_{n-2}(X)
\]
に対し $c_N(x),c_N(y),c_N(z),\Omega_q$ が全中心を生成する。
特に非捩れまたは $N>3$ の場合、この固定包絡上の全H3候補は
$\Omega_{q,h}+Dt^3$ のpencilに限られる。
全候補のrank-one整域条件は前定理で満たされており、
残る問題はtailsの有限大域次元の正確な判定である。
\end{proposition}
\begin{proof}
非捩れでは $Z(S_q)=k[xyz]$ であり、最高項の帰納法を用いる。
有限位数では $Z(S_q)$ は $x^N,y^N,z^N,xyz$ で生成される。
全てが持ち上がることを示せば同じ帰納法で十分である。

右側の $x$ を変数 $r$ と書けば、左からの $x$ の作用は
\[
 x\begin{pmatrix}y&z\end{pmatrix}
 =\begin{pmatrix}y&z\end{pmatrix}
 M(r),\qquad
 M(r)=\begin{pmatrix}qr&-q^{-1}\\1&q^{-1}r\end{pmatrix}
\]
となる。
$r=\kappa(u+u^{-1})$ を代入すると二つの固有値は
$\kappa(qu+q^{-1}u^{-1})$ と $\kappa(q^{-1}u+qu^{-1})$ である。
$q^N=1$ だから、両者に $c_N$ を作用させた値は $c_N(r)$ に等しい。
従って $c_N(M(r))=c_N(r)I$ が多項式恒等式として成立し、
$c_N(x)$ は $y,z$ と可換する。循環的に他も従う。

非捩れまたは $N>3$ では中心三次は
$k\Omega_{q,h}\oplus kt^3$ である。
CY包絡ではH3三次は中心であり、$t^3$ のみでは商が整域でないので除外する。
$R_{q,D}$ は整域でHilbert値も三次曲面型だから、
Gabriel--Krull critical性とrank-one torsion-free条件を満たす。
tailsが有限大域次元なら、AS Gorenstein parameter1の双対・飽和・AZ認識によりH3を得る。
一般 $q$ でこの有限大域次元条件が未判定であることは保持する。
\end{proof}

差分係数の分子が重根となる二つの値は
\[
 D_+=-\frac{q}{(q-1)(q+1)^3},\qquad
 D_-=-\frac{q}{(q-1)^3(q+1)}
\]
である。原始三乗根では、既に完全判定した $f,-f/3$ に一致する。
しかし一般 $q$ でこの二値だけが非H3であるとはまだ証明していない。
分母の特殊軌道、有限次元表現とtailsの局所条件をさらに調べる必要がある。
\texttt{verify\_\allowbreak cubic\_\allowbreak quantum\_\allowbreak triangle\_\allowbreak difference.py} は、
三関係、Casimir値、分子判別式を全 $q,D$ の有理関数恒等式として検算する。

\begin{proposition}[全候補pencilの幾何学的矢]
\label{prop:quantum-triangle-general-arrows}
$\mathcal Q=\{(q,D):q\ne0,\pm1\}$ とする。
$R_{q,D}$ の全graded同型を次数スカラーで割ると、同型群oidは
\[
 [\mathcal Q/\mathfrak S_4]
\]
となる。ここで $\mathfrak A_4$ は $\mathcal Q$ 上自明に作用し、奇置換は
\[
 (q,D)\longmapsto(q^{-1},-q^2D)
\]
として作用する。従って全候補表示のrigidified stackは分離二次元DM stackで、慣性は常に $\mathfrak A_4$ である。
H3部分へ制限すれば全intrinsic pointed矢を与えるが、一般 $q$ での許容部分の開性はまだ主張しない。
\end{proposition}
\begin{proof}
中央一次直線は $kt$ に限るので、同型は $t\mapsto dt'$、$X\mapsto MX'+vt'$ と書ける。
先頭の三量子関係の平方係数を比較すると
$(1-q)M_{ik}M_{jk}=0$、$i\ne j$ となる。
局所底上で $M$ は可逆だから、各列に単元が存在し、他の成分は零となる。
従って非被約底でも $M$ はmonomialである。
循環置換では $q'=q$、奇置換では $q'=q^{-1}$ に限る。
$vt'$ は非対角の一次係数比較で消える。

循環置換の場合、共通スカラーを除く自由度は二符号変更と循環置換だけである。
奇置換の代表 $x\leftrightarrow y,z\mapsto z$ は $t\mapsto-qt'$ と組み合わせれば
$q$ の関係を $q^{-1}$ の関係に移す。
このとき $\Omega_{q,h}\mapsto q\Omega_{q^{-1},h}$ なので
$D'=-q^2D$ となる。
二度作用すれば恒等に戻り、二符号変更と全三置換は
$(\mu_2)^2\rtimes\mathfrak S_3\simeq\mathfrak S_4$ をなす。
$q\ne\pm1$ では奇置換は基底点を固定しないから慣性は正確に $\mathfrak A_4$ である。
平方係数比較は全基底で成り立ち、有限群作用によるquotientなのでstackの主張も従う。
\end{proof}

\subsubsection{奇数位数の全中心と特異切断の有限一覧}
$q$ の正確な位数を奇数 $N\ge3$ とし、$h=q^2-1$ と置く。
$P_0(W)=2,P_1(W)=W,P_n(W)=WP_{n-1}(W)-P_{n-2}(W)$ とする。
$q$-so3中心のChebyshev表示は既知である\cite{HuangAWCentre}。
ここでは直前の忠実模型から、本稿のCasimir正規化で直接導く。
新規性を主張するものではない。

\begin{theorem}[奇数位数のFricke中心と二つの非許容端点]
\label{thm:quantum-triangle-odd-centre}
$X=c_N(x),Y=c_N(y),Z=c_N(z)$、
\[
 w_D=-(q+q^{-1})-\frac{h^3}{q^2}D,\qquad
 \alpha=h^{-N},\qquad
 \beta_N(D)=-h^{-3N}\bigl(2+P_N(w_D)\bigr)
\]
と置く。このとき
\[
 Z(U_{q,D})=
 k[X,Y,Z]/\bigl(XYZ+\alpha(X^2+Y^2+Z^2)+\beta_N(D)\bigr).
\]
$U_{q,D}$ のPI次数は $N$ である。
中心が特異な $D$ は正確に $N+1$ 個で、
$P_N(w_D)=-2$ なら原点に一つの $A_1$、
$P_N(w_D)=2$ なら四点 $-2\alpha(\epsilon_x,\epsilon_y,\epsilon_z)$
（$\epsilon_x\epsilon_y\epsilon_z=1$）に $A_1$ を持つ。

$w_D=-2,2$ に対応する値はそれぞれ上記 $D_+,D_-$ であり、
この二つの切断はH3でない。
残る $N-1$ 個の特異中心値を持つ切断については、
この定理だけでH3とも非H3とも結論しない。
\end{theorem}
\begin{proof}
差分環内で $a=u^N,s=\mathcal T^N$、$K=\kappa^N$、$\nu=\rho^N$ と書く。
$N$ は奇数だから $\nu^2=1$、$K^2=h^{-2N}$、$\nu/K=-h^N$ である。
$X=K(a+a^{-1})$ が成立する。
$Y$ は $x$ と可換するので、その $\mathcal T$ 指数は $N$ の倍数である。
一方、奇数多項式 $c_N(y)$ の指数は奇数で絶対値が $N$ 以下だから、
$\pm N$ の二項だけが残る。同様に $Z$ も計算でき、
\[
 Y=s+P/s,\qquad Z=\nu(as+a^{-1}P/s),\qquad
 P=\prod_{j=0}^{N-1}b_D(q^ju)
\]
となる。
$b_D$ の分子を $q(v-r_1)(v-r_2)$、$v=u^2$ と書くと
$r_1r_2=q^2$、$(r_1+r_2)/q=w_D$ である。
$q^2$ も位数 $N$ なので積を厳密に取れ、
\[
 P=\frac{a^4-P_N(w_D)a^2+1}{h^{2N}(a^2-1)^2}
\]
を得る。
$Y,Z$ から $s$ を消去すれば
\[
 Z^2-\frac{\nu X}{K}YZ+Y^2+
 P\bigl((X/K)^2-4\bigr)=0
\]
となり、表示したFricke関係が従う。

$D=-\Omega_q$ を独立変数として同じ恒等式を用いると、
$Z(V_q)$ の全関係はこの一本である。
実際、$k[X,Y,Z]$ 上の $1,\Omega_q,\ldots,\Omega_q^{N-1}$ は
最高PBW記号の指数を $N$ で割った余りから独立であり、
関係は $\Omega_q$ に関して次数 $N$ のモニック式である。
また $V_q$ 自身は $k[X,Y,Z]$ 上PBW基底 $x^iy^jz^k$、$0\le i,j,k<N$ を持つ。
従って全中心上のgeneric階数は $N^2$ である。

固定 $D$ のFricke環を $C_D$ とする。
偏微分を解くと特異点は主張の有限集合だけなので、$C_D$ は正規整域である。
$U_{q,D}$ は $C_D$ 上有限であり、filtered Hilbert多項式の最高係数の比は $N^2$。
その中心分数体に余分な拡大がないことも次のように確認できる。
$L=k(a,s)$ はFricke分数体の二次拡大で、
$a$ は $a+a^{-1}=X/K$ を満たす。
$x/\kappa=u+u^{-1}$ と $u^N=a$ の二式から $u$ を有理的に復元でき、
さらに $y,z$ から $\mathcal T$ を復元できる。
従ってgeneric代数を $L$ へ拡大すると位数 $N$ の差分巡回代数そのものとなる。
両辺の階数は $N^2$ なので同型であり、generic中心は正確に $\operatorname{Frac}C_D$ である。
$C_D$ の正規性と中心元の整性から $Z(U_{q,D})=C_D$ を得る。

Fricke偏微分の解は原点または上記四点に限る。
それぞれのHessian行列式は $8\alpha^3,-32\alpha^3$ で非零である。
$P_N(W)\mp2$ の根は、端点 $\pm2$ だけ単根で、他は重根となる。
両式の根の和集合は $N+1$ 個であり、$w_D$ は $D$ の可逆な一次式だから個数も従う。

最後に、$V_q$ は上の多項式部分環上自由、中心上MCMである。
そのCY双対・filtered最大整環性と正則中心上のAzumaya性の証明は、
三乗根の場合の証明で指数3を $N$ に置き換えればそのまま成立する。
具体的には先頭の量子多項式環がtoric最大整環であり、
$\operatorname{Hom}_{Z(V_q)}(V_q,Z(V_q))\simeq V_q$、
余次元1非分岐性、正則点でのMCM自由性を順に用いる。
端点では $P_N'(\pm2)\ne0$ だから、切断中心の全二重点で
三次元中心 $Z(V_q)$ は正則である。
従って完備局所化した切断は特異な $A_1$ 中心上の行列環となり、
大域次元は無限である。affine開集合で既に失敗するのでH3でない。
内部の重根では三次元中心自身が通常二重点となり、
このAzumaya除外を適用できないことに注意する。
\end{proof}

\begin{lemma}[奇数位数の境界orderとaffine判定]
\label{lem:quantum-triangle-odd-boundary}
前定理の奇数位数で、$R_{q,D}^{(N)}$ の全中心のProjは
\[
 XYZ+\alpha T(X^2+Y^2+Z^2)+\beta_N(D)T^3=0
 \quad\subset\mathbb P^3_{T,X,Y,Z}
\]
である。対応するorderは無限遠三角形に沿ってterminalである。
従ってH3判定はaffine $U_{q,D}$ の有限大域次元と同値である。
中心が滑らかな切断は全てH3となる。
中心の $A_1$ 点では、次元が $1,\ldots,N-1$ の既約表現が少なくとも一つ存在することが、
局所有限大域次元の必要十分条件である。
\end{lemma}
\begin{proof}
$B_q$ は $k[t,X_h,Y_h,Z_h]$ 上自由階数 $N^3$ である。
その正則中心三次商は中心hypersurface上MCMとなる。
重み $(1,N,N,N)$ の中心hypersurfaceは、affineの孤立二重点と
無限遠の三本の軸以外では正則なので正規である。
前定理のgeneric中心と整性により全中心と一致する。
$N$-Veroneseを取れば表示した通常の三次曲面となる。

無限遠三角形上でこの中心曲面は滑らかである。
例えば辺 $X=0$ で $Y,Z$ を可逆にすれば $y,z$ が単元となる。
$c_N$ の斉次化を
$Y_h=y^N+\sum_{j\ge1}a_jt^{2j}y^{N-2j}$ と書き、$v=ty^{-1}$ と置くと
\[
 v^N=\frac{T}{Y_h}\left(1+\sum_{j\ge1}a_jv^{2j}\right),
 \qquad J=(v),\quad J^N=(T).
\]
residueは $u=zy^{-1}$ による $k(u)/k(u^N)$ である。
三乗根の場合と同じ主radicalの証明でgeneric辺orderは非可換DVRとなり、
coverは両頂点で完全分岐指数 $N$ を持つ。
三辺全てに同じ証明を適用し、terminal判定と局所模型
\cite[Theorems 2.7, 2.17]{CHIOrders}を用いる。

affineでは $U_{q,D}$ は正規中心上の対称MCM非分岐最大整環である。
これは $V_q$ の正則中心上のAzumaya性を余次元1で制限し、
中心商の双対と反射性を用いれば得られる。
$A_1$ の被覆群は $\mu_2$ であり、完備局所orderの二つのMcKay重複度は
$r,s\ge0$、$r+s=N$ を満たす。
両方正なら局所大域次元2、一方零なら特異中心上の行列環となる。
従って次元 $0<r<N$ のsimpleが一つ見つかれば両方正となる。
逆も同じ重複度記述から従う。
最後にAS Gorenstein parameter1、整域性、AZ飽和とSerre認識を適用する。
\end{proof}

\begin{theorem}[位数5三角境界pencilの全H3切断]
\label{thm:quantum-triangle-order5}
$q$ が原始五乗根なら
\[
 R_{q,D}\text{ が強いH3}\quad\Longleftrightarrow\quad
 (D-D_+)(D-D_-)\ne0.
\]
四つの内部特異中心値も全てH3である。
このpencil内の全intrinsic幾何学的矢は
命題\ref{prop:quantum-triangle-general-arrows}の制限で与えられる。
\end{theorem}
\begin{proof}
$\theta=q+q^{-1}$、$\eta=q^2+q^{-2}$ と置く。
端点以外の四つの $w_D$ は $-\theta,-\eta,\theta,\eta$ で相異なる。
負の二値では原点に一つの $A_1$、正の二値では四つの $A_1$ がある。
$w_D=-\theta$ は $D=0$ で零characterを持つ。
$w_D=\eta$ は非零四character
$(x,y,z)=(\epsilon_x,\epsilon_y,\epsilon_z)/(1-q)$ に対応する。

残る二値には次の二次元既約表現を用いる。
これらの関係恒等式自体は一般 $q$ の有理式で成立する。
まず $A=\sqrt{-1}/(q+1)$ として
\[
 x=\begin{pmatrix}A&0\\0&-A\end{pmatrix},\quad
 y=\begin{pmatrix}0&1\\A^2&0\end{pmatrix},\quad
 z=\begin{pmatrix}0&(1+q)A\\-(1+q)A^3&0\end{pmatrix}
\]
は $w_D=-(q^2+q^{-2})=-\eta$ を持つ。
次に
\[
 a=\frac1{1-q},\quad b=a(q+q^{-1}-1),\quad
 c=\frac b{q-1},\quad d=\frac{q^2+1}{q(q-1)}
\]
として
\[
 x=\begin{pmatrix}a&0\\0&b\end{pmatrix},\quad
 y=\begin{pmatrix}d&1\\c&0\end{pmatrix},\quad
 z=\begin{pmatrix}d&q\\q^{-1}c&0\end{pmatrix}
\]
は $w_D=q^4+q^{-4}=\theta$ を持つ。
両者で $x$ の固有値は相異なり、$y$ の両非対角成分は非零なので既約である。
前者の中心値は原点、後者の中心三座標は全て非零である。
後者に二符号変更群を作用させれば四つの $A_1$ 全てを覆う。
これらは $\Phi_5(q)$ による厳密な係数還元でも確認している。

全内部二重点に次元1または2のsimpleが存在し、いずれもPI次数5より小さい。
前補題により全て局所有限大域次元となる。
端点は前定理で除外済み、他の切断は滑らかな中心を持つので主張が従う。
\end{proof}

\texttt{verify\_\allowbreak cubic\_\allowbreak quantum\_\allowbreak triangle\_\allowbreak \allowbreak odd\_\allowbreak center.py} は
消去恒等式、位数3,5,7,9の巡回積、三乗根での独立PBW中心式との一致、
Fricke Hessianを厳密検算する。
\texttt{verify\_\allowbreak cubic\_\allowbreak quantum\_\allowbreak triangle\_\allowbreak \allowbreak order5.py} は
全有理式での二次元表現恒等式と、位数5での中心値・既約性・四内部値の区別を検算する。
次に、この二次元表示を任意次元の三重対角表示へ拡張して、全奇数位数を扱う。

\begin{theorem}[全奇数位数の三角境界pencilの完全判定]
\label{thm:quantum-triangle-all-odd}
$q$ が任意の奇数位数 $N\ge3$ を持つとき、
\[
 R_{q,D}\text{ が強いH3}\quad\Longleftrightarrow\quad
 (D-D_+)(D-D_-)\ne0.
\]
全 $N-1$ 個の内部特異中心切断もH3である。
ただし、この結論は偶数位数または非捩れ $q$ の許容条件を含まない。
\end{theorem}
\begin{proof}
$m=(N-1)/2$ とする。各 $1\le d\le m$ に対して次元 $d$ の既約表現を二種類構成する。
$d=1$ は既述の零characterと四つの非零characterである。
以下 $d\ge2$ とし、添字は $j=0,\ldots,d-1$ とする。
$h=q^2-1$ と置く。

第一の型では $R^2=-q^{2-d}$ を選び、
\[
 \lambda_j=-\frac{Rq^j+q^{1-j}R^{-1}}h,\qquad
 A_j=Rq^{j-1}-q^{2-j}R^{-1},
\]
\[
 c_j=\frac{q^2\{q^d+q^{-d}-q^{2j-d}-q^{d-2j}\}}
                  {h^2A_jA_{j+1}}\quad(1\le j<d),\qquad c_0=c_d=0
\]
と置く。$x$ は対角 $\lambda_j$、$y$ は対角零、上隣接成分1、下隣接成分 $c_j$ とする。
$z=xy-qyx$ は対角零で、上隣接成分 $Rq^j$、下隣接成分 $R^{-1}q^{1-j}c_j$ となる。
三関係のうち二つは定義と $\lambda_j$ の式から従い、残りは
\[
 c_jA_j-c_{j+1}A_{j+2}=\lambda_j
\]
である。これは等比級数を和すれば成立する。
両端では $c_0=c_d=0$ である。
最後の対角成分でCasimirを評価すると
\[
 w_D=-(q^d+q^{-d})
\]
を得る。
$A_j=0$ や異なる $\lambda_j$ の一致は $q$ の冪が $-1$ になることを要求するので起こらない。
$c_j=0$ は $q^{2j-d}=q^{\pm d}$ を要求するが、$0<j<d\le m$ では起こらない。
従って両向きに全隣接成分が非零で、対角 $x$ の固有値は相異なる。
補間多項式で各固有空間を射影できるため表現は既約である。

第二の型では
\[
 \lambda_j=-\frac{q^{j+1}+q^{-j}}h,\qquad
 A_j=q^j-q^{1-j},\qquad
 c_j=\frac{q^2\{q^{2j}+q^{-2j}-q^{2d}-q^{-2d}\}}
                  {h^2A_jA_{j+1}}\quad(1\le j<d)
\]
と置く。$x$ は対角 $\lambda_j$、$y$ の上隣接成分は1、下隣接成分は $c_j$ とし、
対角は $j=0$ のみ
\[
 d_0=\frac{q^d-q^{-d}}{(q-q^{-1})(q-1)}
\]
とする。$z=xy-qyx$ の上隣接成分は $q^{j+1}$、下隣接成分は $q^{-j}c_j$、
対角は $j=0$ のみ $d_0$ となる。
内部の関係は同じ再帰式で、最初の対角式だけ
\[
 (1-q)d_0^2-c_1A_2=\lambda_0
\]
となる。最後は $c_d=0$ であり、全て表示式で満たされる。
Casimir値は
\[
 w_D=q^{2d}+q^{-2d}
\]
である。
分母には $q^{2j-1}-1$ と $q^{2j+1}-1$ が現れるが、
$1\le j<d\le m$ ではその指数は $1,\ldots,N-2$ に収まる。
$c_j=0$ なら $N\mid2(d-j)$ または $N\mid2(d+j)$ となり不可能である。
異なる $\lambda_j$ の一致も $N\mid j+k+1$ を要求し、
$1\le j+k+1\le2d-1<N$ なので起こらない。
従ってこの型も既約である。

$-(q^d+q^{-d})$、$d=1,\ldots,m$ は、$P_N(w)=-2$ の端点以外の全根を相異なって走る。
各々の唯一の $A_1$ 点に次元 $d<N$ のsimpleを得る。
同様に $q^{2d}+q^{-2d}$ は $P_N(w)=2$ の端点以外の全根を走る。
これは $N$ が奇数で2が可逆であることによる。
この場合中心の四点に二符号変更群を作用させれば全てに次元 $d<N$ のsimpleを得る。
前補題のMcKay判定とterminal境界から、内部特異値も全てH3となる。
端点二値の除外と滑らかな中心の許容性は既証である。
\end{proof}

\texttt{verify\_\allowbreak cubic\_\allowbreak quantum\_\allowbreak triangle\_\allowbreak \allowbreak odd\_\allowbreak representations.py} は
$Q=q^d,J=q^j$ を形式変数として、両型の全再帰式・端点式とCasimir式を厳密検算する。
古典型では $R^2=-q^2/Q$ での多項式剰余を取り、有限個の位数での標本検算へ還元していない。
分母の非消滅・既約性・全中心点の被覆は上記の指数範囲による証明である。
